% =============================================================================
%                             PROBABILITY PLUGIN
% =============================================================================
%
% Probability measures, expectations, distributions, covariance, and common
% statistical operators.
%
% Requires:
%   xamsmath core helpers and mathcommand PIE support.

% =============================================================================
%                         CONDITIONAL DELIMITERS
% =============================================================================

% -----------------------------------------------------------------------------
% Internal Commands
% -----------------------------------------------------------------------------
%
% \BracesSplit
%
% Purpose:
%   Parenthesized split delimiter for conditional probability-like expressions.
%
% Parameters:
%   #1 - Expression before a vertical bar.
%   #2 - Expression after a vertical bar.
%
% Usage:
%   Internal use by probability functionals.

\DeclarePairedSplitDelimiter{\BracesSplit}{(}{)}

% -----------------------------------------------------------------------------
% Internal Commands
% -----------------------------------------------------------------------------
%
% \BraketSplit
%
% Purpose:
%   Bracketed split delimiter for expectation-like expressions.
%
% Parameters:
%   #1 - Expression before a vertical bar.
%   #2 - Expression after a vertical bar.
%
% Usage:
%   Internal use by expectation and moment operators.

\DeclarePairedSplitDelimiter{\BraketSplit}{[}{]}

% =============================================================================
%                         CONDITIONAL TEMPLATES
% =============================================================================

% -----------------------------------------------------------------------------
% Internal Commands
% -----------------------------------------------------------------------------
%
% \ConditionalFunctional
%
% Purpose:
%   Defines a PIE-aware probability-like functional whose optional conditional
%   argument is written in parentheses.
%
% Parameters:
%   #1 - Command to define.
%   #2 - Printed functional head, such as \mathbb{P}.
%
% Generated command parameters:
%   #1 - Prime tokens absorbed by mathcommand.
%   #2 - Index tokens absorbed by mathcommand.
%   #3 - Exponent tokens absorbed by mathcommand.
%   #4 - Optional parenthesized conditional expression.
%
% Usage:
%   \ConditionalFunctional{\P}{\mathbb{P}}
%   \P(A | B)

\NewDocumentCommand{\ConditionalFunctional}{mm}{
    \NewDocumentCommand{#1}{mmm d()}{
        \IfNoValueTF{##4}{#2##1##2##3}{#2##1##2##3\BracesSplit{##4}}
    }
    \piefication{#1}
}

% -----------------------------------------------------------------------------
% Internal Commands
% -----------------------------------------------------------------------------
%
% \ConditionalFunctionalOperator
%
% Purpose:
%   Defines a PIE-aware expectation-like operator whose optional conditional
%   argument is written in square brackets.
%
% Parameters:
%   #1 - Command to define.
%   #2 - Printed operator head, such as \mathbb{E}.
%
% Generated command parameters:
%   #1 - Prime tokens absorbed by mathcommand.
%   #2 - Index tokens absorbed by mathcommand.
%   #3 - Exponent tokens absorbed by mathcommand.
%   #4 - Optional bracketed conditional expression.
%
% Usage:
%   \ConditionalFunctionalOperator{\E}{\mathbb{E}}
%   \E[X | Y]

\NewDocumentCommand{\ConditionalFunctionalOperator}{mm}{
    \NewDocumentCommand{#1}{mmm d[]}{
        \IfNoValueTF{##4}{#2##1##2##3}{#2##1##2##3\BraketSplit{##4}}
    }
    \piefication{#1}
}

% =============================================================================
%                         PROBABILITY FUNCTIONALS
% =============================================================================

% -----------------------------------------------------------------------------
% User Commands
% -----------------------------------------------------------------------------
%
% \P
%
% Purpose:
%   Typesets a probability measure with optional PIE notation and conditional
%   parentheses.
%
% Usage:
%   \P(A)
%   \P_X(A | B)

\let\P\relax
\ConditionalFunctional{\P}{\mathbb{P}}

% -----------------------------------------------------------------------------
% User Commands
% -----------------------------------------------------------------------------
%
% \Law
%
% Purpose:
%   Typesets the law or distribution of a random variable, with optional
%   conditional parentheses.
%
% Usage:
%   \Law(X)
%   \Law(X | Y)

\ConditionalFunctional{\Law}{\mathrm{Law}}

% =============================================================================
%                         EXPECTATION AND MOMENTS
% =============================================================================

% -----------------------------------------------------------------------------
% User Commands
% -----------------------------------------------------------------------------
%
% \E
%
% Purpose:
%   Typesets expectation with optional PIE notation and conditional brackets.
%
% Usage:
%   \E[X]
%   \E_{\P}[X | Y]

\ConditionalFunctionalOperator{\E}{\mathbb{E}}

% -----------------------------------------------------------------------------
% User Commands
% -----------------------------------------------------------------------------
%
% \D
%
% Purpose:
%   Typesets variance or dispersion notation using blackboard-bold D.
%
% Usage:
%   \D[X]

\ConditionalFunctionalOperator{\D}{\mathbb{D}}

% -----------------------------------------------------------------------------
% User Commands
% -----------------------------------------------------------------------------
%
% \std
%
% Purpose:
%   Typesets standard deviation as an upright statistical operator.
%
% Usage:
%   \std[X]

\ConditionalFunctionalOperator{\std}{\mathrm{std}}

% -----------------------------------------------------------------------------
% User Commands
% -----------------------------------------------------------------------------
%
% \mad
%
% Purpose:
%   Typesets mean absolute deviation as an upright statistical operator.
%
% Usage:
%   \mad[X]

\ConditionalFunctionalOperator{\mad}{\mathrm{mad}}

% -----------------------------------------------------------------------------
% User Commands
% -----------------------------------------------------------------------------
%
% \Skew
%
% Purpose:
%   Typesets skewness as an upright statistical operator.
%
% Usage:
%   \Skew[X]

\ConditionalFunctionalOperator{\Skew}{\mathrm{skew}}

% -----------------------------------------------------------------------------
% User Commands
% -----------------------------------------------------------------------------
%
% \Kurt
%
% Purpose:
%   Typesets kurtosis as an upright statistical operator.
%
% Usage:
%   \Kurt[X]

\ConditionalFunctionalOperator{\Kurt}{\mathrm{kurt}}

% =============================================================================
%                                COVARIANCE
% =============================================================================

% -----------------------------------------------------------------------------
% Internal Commands
% -----------------------------------------------------------------------------
%
% \covbraces
%
% Purpose:
%   Parenthesized two-argument delimiter used by \Cov.
%
% Parameters:
%   #1 - First covariance argument.
%   #2 - Second covariance argument.
%
% Usage:
%   Internal use by \Cov.

\DeclarePairedDelimiterX\covbraces[2]{(}{)}{#1,#2}

% -----------------------------------------------------------------------------
% User Commands
% -----------------------------------------------------------------------------
%
% \Cov
%
% Purpose:
%   Typesets covariance with parenthesized arguments.
%
% Parameters:
%   #1 - First random variable or expression.
%   #2 - Second random variable or expression.
%
% Usage:
%   \Cov{X}{Y}

\newmathcommand{\Cov}[2]{\mathrm{Cov}\covbraces{#1}{#2}}

% =============================================================================
%                               DISTRIBUTIONS
% =============================================================================

% -----------------------------------------------------------------------------
% User Commands
% -----------------------------------------------------------------------------
%
% \Uniform
%
% Purpose:
%   Typesets the uniform distribution symbol.
%
% Usage:
%   X \sim \Uniform(a,b)

\newmathcommand{\Uniform}{\mathcal{U}}

% -----------------------------------------------------------------------------
% User Commands
% -----------------------------------------------------------------------------
%
% \Normal
%
% Purpose:
%   Typesets the normal distribution symbol.
%
% Usage:
%   X \sim \Normal(\mu,\sigma^2)

\newmathcommand{\Normal}{\mathcal{N}}

% -----------------------------------------------------------------------------
% User Commands
% -----------------------------------------------------------------------------
%
% \Poiss
%
% Purpose:
%   Typesets the Poisson distribution name.
%
% Usage:
%   X \sim \Poiss(\lambda)

\newmathcommand{\Poiss}{\mathrm{Poiss}}

% -----------------------------------------------------------------------------
% User Commands
% -----------------------------------------------------------------------------
%
% \Exp
%
% Purpose:
%   Typesets the exponential distribution name.
%
% Usage:
%   X \sim \Exp(\lambda)

\newmathcommand{\Exp}{\mathrm{Exp}}

% =============================================================================
%                            STOCHASTIC PROCESS TEXT
% =============================================================================

% -----------------------------------------------------------------------------
% User Commands
% -----------------------------------------------------------------------------
%
% \cadlag
%
% Purpose:
%   Typesets the French term for right-continuous paths with left limits.
%
% Usage:
%   \cadlag

\newtextcommand{\cadlag}{c\`adl\`ag}

% =============================================================================
%                                USAGE NOTES
% =============================================================================
%
% Probability measures:
%   \P(A)
%   \P(A | B)
%   \P_X(A | B)
%
% Expectation and moments:
%   \E[X]
%   \E_{\P}[X | Y]
%   \D[X]
%   \std[X]
%
% Statistics and distributions:
%   \Cov{X}{Y}
%   X \sim \Normal(\mu,\sigma^2)
%   X \sim \Uniform(a,b)
